\everymath{\displaystyle}

%%%%%%%%%%%%%%%%%%%%%%%%%%%%%%%%%%%%%%%%
%           main body
%%%%%%%%%%%%%%%%%%%%%%%%%%%%%%%%%%%%%%%%

%-----------------------------------
%	TITLE PAGE
%-----------------------------------

\title[第八章\\
反常积分]{\texorpdfstring{\LARGE \bfseries 第八章\quad 反常积分}{第八章 反常积分}} % The short title appears at the bottom of every slide, the full title is only on the title page
\author[\smaller \S 8.1\\
反常积分概念和计算] % Your institution as it will appear on the bottom of every slide, maybe shorthand to save space
{\Large \bfseries \S 8.1 反常积分的概念和计算}
% \author{Singular Value Decomposition} % Your name
% \date{\today} % Date, can be changed to a custom date
\date{}

%------------------------------------------------

\begin{document}

%------------------------------------------------

\begin{frame}

\titlepage % Print the title page as the first slide

\end{frame}

%------------------------------------------------

% \begin{frame}
% \frametitle{内容提要} % Table of contents slide, comment this block out to remove it
% \tableofcontents % Throughout your presentation, if you choose to use \section{} and \subsection{} commands, these will automatically be printed on this slide as an overview of your presentation
% \end{frame}

%------------------------------------------------

\begin{frame}

\titlepage % Print the title page as the first slide

\end{frame}

%-----------------------------------
%	PRESENTATION SLIDES
%-----------------------------------

%------------------------------------------------

\begin{frame}
\frametitle{绪言}

\textbf{问题:} 之前利用函数显式表达 $y = \sqrt{1 - x^2}$ 计算半圆弧长的时候, 遇到了如下的积分:
\[
\ell = \int_{-1}^1 \sqrt{1 - (y')^2} ~ \mathrm{d} x = \int_{-1}^1 \sqrt{1 - \left(-\frac{x}{\sqrt{1 - x^2}}\right)^2} ~ \mathrm{d} x = \int_{-1}^1 \frac{1}{\sqrt{1 - x^2}} ~ \mathrm{d} x.
\]
被积函数在积分区间的端点处无定义, 且在该区间上无界, 这是黎曼积分的无法处理的. 但这并不表明弧长是无法计算的. 所以我们需要拓展黎曼积分的定义.

\pause
\vspace{0.5em}

一个自然的想法是, \alert{``取黎曼积分的极限''}, 对于上述的积分, 就是
\[
\lim_{\varepsilon \to 0^+} \int_{-1 + \varepsilon}^0 \frac{1}{\sqrt{1 - x^2}} ~ \mathrm{d} x
~~ \text{与} ~~
\lim_{\varepsilon \to 0^+} \int_0^{1 - \varepsilon} \frac{1}{\sqrt{1 - x^2}} ~ \mathrm{d} x.
\]

\end{frame}

%------------------------------------------------

\section{反常积分的定义}

%------------------------------------------------

\begin{frame}
\frametitle{反常积分的定义}


\end{frame}

%------------------------------------------------

\end{document}
